\documentclass[2pt]{book}
\usepackage{graphicx}
\usepackage{html}
\usepackage{hyperref}
\usepackage{lmodern}
\usepackage{tocloft}
\usepackage{svg}
\setlength{\cftpartindent}{0em}
\setlength{\cftchapindent}{0em}
\setlength{\cftsecindent}{0em}
\setlength{\cftsubsecindent}{0em}
\setlength{\cftsubsubsecindent}{0em}
\graphicspath{{Images/}}
\author{Doug Royer / RiverExplorer™}
\date{13-MAR-2025}
\title{RPCGEN++ An XDR and ABNF compiler and code generator}
\begin{htmlonly}
  <link rel='stylesheet' type='text/css' href='rpcgen++.css'/>
\end{htmlonly}
%
% Define Abstract ..
%
\newenvironment{abstract}
{%
  \small
  \begin{center}
    \bfseries{abstract}
  \end{center}
}
{}

\begin{document}
\maketitle
\frontmatter
\begin{figure}
  \centering
  \includegraphics{RiverExplorerLLC-Logo-BLACK-131x50-Transparent.ps.png}
\end{figure}

\begin{abstract}

  \includegraphics{UnderConstruction.png}
  \underline{\textbf{THIS IS WORK IN PROGRESS}}
  
  RPCGEN++ is an Open Source XDR and ABNF processor and code generator.
  RPCGEN++ can read \href{https://ietf.org}{IETF} data description
  formats;
  \href{https://datatracker.ietf.org/doc/html/rfc4506}{[RFC-4506]}
  XDR and \href{https://datatracker.ietf.org/doc/html/rfc5234}{[RFC-5234]}
  ABNF input files and produce runnable code for Linux, Windows,
  and Android. The resulting code has also been tested on the Android
  Oculus virtual reality headset.

  The code generated can connect systems across the internet,
  or your intranet.
  RPCGEN++ takes care of encoding, transferring, decoding and over
  TLS. This RPCGEN++ processor can also create the client and server side
  prototypes where you fill in the details of your design.
  RPCGEN++ also allows for describing bit field data for XDR and ABNF.
  
  The XDR side of the implementation was inspired by the ONC
  rpcgen program written by Sun Microsystems
  in the 1990's.
  This version has extensions that allow for C++ like, classes, class
  methods, and comments to be in included into the (.x) source XDR
  description files. RPCGEN++ was written using the antlr4 protocol
  compiler. It produces C++ code, and has been designed to be able
  to create other languages.

  The ABNF code is from scratch and allows the user to annotate
  the ABNF with instructions about the intent of the structure
  of the intended program.
  
  RPCGEN++ was originally written to support the Phoenix project.
  Its source is part of the Phoenix project.
  (\href{https://github.com/RiverExplorer/Phoenix}{Phoenix Open Source Project})
  RPCGEN++ and Phoenix are both Open Source.
\end{abstract}
\tableofcontents
\listoffigures
\listoftables

\mainmatter
\chapter{Overview}
\section{Introduction}
The client and the server send data to the network, and get
data from the network.
This \mbox{input/output} model (IO Model) works with
both the Phoenix client and Phoenix server.

This model support both threaded and non threaded
implementations.
For non-threaded systems the data is processed serially
or however the implementation handles both incoming
and outgoing data.
The data flow is the same for both threaded and non threaded
implementations.

\begin{figure}
  \centering
  \caption{I/O Model}
  \label{fig:IOModel}
  \includesvg{Generated/IO-00.svg}
\end{figure}

\begin{description}
\item[(A)Main Thread]\label{item:MainThread}
	The main thread processes data and formats it into
  IOVectorPacket objects that contains the command
  and references to the data that will be sent.
  Then stores the IOVectorPackt (figures: \ref{fig:IovPacket} and \ref{fig:IovPacket2})
  into queue \textit{OQ}.
  The client \textit{Main Thread} is described in section \ref{sec:ClientMainThread}
  And the server \textit{Main Thread} is described in section \ref{sec:ServerMainThread}

\item[(B) Output Thread]
  This thread waits for the output queue \textit{OQ} to have
  content.
  It checks to make sure that the network is writable and not write blocked
  or otherwise busy.
  When it can proceed, it takes an entry out of the queue,
  and (C) XDR encodes it.
  The XDR stream is opened to the TLS enctyption.
  The \textit{Output Thread} is described in section \ref{sec:OutputThread}
  
\item[(C) XDR Encoding]\label{item:XdrEncodingThread}
	This thread, takes one packet at at time
	out of the input queue,	converts it from the native computer
  to network compatible XDR format.
  It reads the IOVectorPacket data to an XDR stream and that streams is sent directly
  to the \textit{(D) TLS Encryption} code.
  \textit{XDR Encoding} is described in section \ref{sec:XdrEncode}

\item[(D) TLS Encryption]\label{item:TlsEncryption}
  This code takes the XDR encoded data from its input stream
  and encrypts it to the network.
  Which encrypts it according to the certificates supplied by the configuration.
  and sends it out to the network.
  \textit{TLS Encryption} is described in section \ref{sec:TlsIOOut}

\item[(E) Input Thread]\label{item:InputThread}
  This thread waits for any data coming from the network
  and feeds it into the \textit{TLS Decryption} code.
    
\item[(F) TLS Decryption]\label{item:TlsDecryptionThread}
  This code decrypts it on the fly, and streams it to \textit{(G) XDR Decoding} code.
  \textit{TLS Decryption} is described in section \ref{sec:TlsIOIn}

\item[(G) XDR DECODING]\label{item:XdrDecoding}
  This takes data from the \textit{TLS Decryption} code and decodes
  it back to native computer data format.
  \textit{XDR Decoding} is described in section \ref{sec:XdrDecode}
  
\item[(H) Dispatch Thread]\label{DispatchThread}
	This is the dispatch thread.
  It takes the decoded packet and dispatches it to the code that can handle the data.
  The \textit{Dispatch Thread} is described in section \ref{sec:DispatchThread}

\end{description}


This model is the core of both the client and server.

\section{XDR Differences from RFC-4506}
\subsection{Place 'struct' Outside of 'union'}
Currently RPCGEN++ does not allow structs to be defined in unions.
Define them out side of the union, and then use them.

For example, this is valid as defined in RFC-4506:
\begin{figure}
\begin{verbatim}
union MasterOfAllTrades (uint32_t Control) {

case 0:
 struct {
   string item<>;
   stringlist Next;
 } element;

 case 1:
   .
   .
   .
\end{verbatim}
\caption{Structure Inside of Union}
\label{fig:StructInFigure}
\end{figure}


Currently, with RPCGEN++, you will need to define your
structure outside of the union as shown in this example:
\begin{figure}
\begin{verbatim}
 struct {
   string item<>;
   stringlist Next;
 } elementObject;

union MasterOfAllTrades (uint32_t Control) {

case 0:
 elementObject element;

 case 1:
   .
   .
   .
\end{verbatim}
\caption{Structure Outside of Union}
\label{fig:StructOutsideFigure}
\end{figure}


\subsection{RPCGEN++ has support for bit XDR fields}
The \textit{bitfield} objects are useful when you have boolean or other
values that do not need to have their own 32-bit storage.
Or they could represent the bit values needed to interface
with hardware.

They are declared by appending a colon (:) and a decimal number
after the data type or variable name.
They are defined using the \textit{bitfield} data type.
In this example, \textit{TwoFourBitValues} is declared
as 8-bits in size, and it contains two 4-bit values.
\begin{figure}
\begin{verbatim}
bitfield TwoFourBitValues:8 {
   int8_t:4 MyVariableName:4;
   int8_t:4 TheOtherFourBits:4;
};
\end{verbatim}
\caption{Two Four Bit values in one 8-bit value}
\label{fig:A32BitExampleXDR}
\end{figure}

The number of bits inside of the \textit{bitfield} must be equal to or
less than the number of bits declared in the bitfield name.

You must declare them with the least significant bits first.
\textit{MyVariableName} in the example would be placed in
bit positions 0-3 and \textit{TheOtherFourBits} would occupy
bit position 4-7.

\textit{For signed values}, make sure that the value you pass
into the encoding call is correctly signed for the bit
width you specified in the XDR file.
If you specify a 4-bit value as signed
and pass it in as an 8-bit signed value, the sign bit will not
be encoded.

The \textit{bitfield} values can be signed or unsigned and
are XDR encoded and XDR decoded on their own.
Once encoded the \textit{bitfield} itself is sent as
opaque data.

In this example, the \textit{bitfield} is 64-bits and
contains multiple values that are not byte aligned:
\begin{figure}
\begin{verbatim}
bitfield ALargeOne:64 {
   int8_t:1 DeviceEnabled:1;     /* 'A' */
   int8_t:2 DeviceMode:2;        /* 'B' */
   int8_t:4 DeviceRotation:4;    /* 'C' */
   int8_t:4 DeviceStatus:4;      /* 'D' */
   int32_t:32 DevicePosition:32; /* 'E' */
};
\end{verbatim}
\caption{Packed 64-bit device example}
\label{fig:ALargeOneXDR}
\end{figure}
The upper 21 bits are undefined in this example,
and are set to zero (0) in this
implementation.
XDR and this implementation sends data in chunks of 32-bits.
The 'A', 'B', 'C', 'D', and 'E' comments in figure
\ref{fig:ALargeOneXDR} are represented as the bits in
the 64-bit figure \ref{fig:ALargeOneBitsXDR}.
\begin{itemize}
\item 'A' is one bit wide.
\item 'B' is two bits wide.
\item 'C' is four bits wide.
\item 'D' is four bits wide.
\item 'E' is 32 bits wide.
\end{itemize}
\begin{figure}
  \includesvg{64bitBitField.svg}
  \caption{64-bit 'ALargeOne' bitfield}
  \label{fig:ALargeOneBitsXDR}
\end{figure}

RPCGEN++ does not impose a size limit on the bit width.
This is an example of a 192-bits object.

\begin{figure}
\begin{verbatim}
bitfield A192BitOne:192 {
  ... define up to 192 bits ...
};
\end{verbatim}
\caption{Packed 192-bit device example}
\label{fig:A192BitXDR}
\end{figure}

\subsection{Added a \textit{class} object that extends \textit{struct}}
RPCGEN++ accepts the keyword \textit{class} to be substituted in
place of \textit{struct}.
When \textit{class} is used, the object may have member functions.
See \ref{ClassDefinition} for the full specification.
\begin{figure}
\begin{verbatim}
class OverTheWirePacket {
  Command    Cmd;
  uint8_t    Payload<>;

  /**
   * You can define a constructor
   *
   * @param WhatToDo the command to process.
   *
   * @param Data In this example, the contents of Data is determinable
   * from WhatToDo.
   */
  ObjectName(Command WhatToDo, uint8_t Data<>);

  /**  
   * When you define a constructor (like the previous line),
   * RPCGEN++ will only create the ones you specify. If you specify none,
   * it will generate a default constructor. So we also define a
   * a default constructor
   */
  ObjectName();

  /**
   * You can define a destructor
   */
  ~ObjectName();

  /**
   * You can define an unlimited number of member functions.
   * The 'const' signifies that the method GetTheAverageValues()
   * does not alter the contents of this object.
   *
   * @param Value1 Any kind or number of function parameters.
   *
   * @param Enabled Another parameter
   */
  float GetTheAverageValues(float Value1, bool Enabled) const;
};

\end{verbatim}
\caption{Example of new class keyword}
\label{fig:ClassIntroXDR}
\end{figure}

\subsection{The \textit{class} data type}\label{sec:ClassDefinition}
RPCGEN++ allows member functions to be placed inside
of class elements that are passed to the target
language as members of the object.

Member names are of the form:
\begin{figure}
\begin{verbatim}
;
; For output languages that do not have pointer or reference types,
; the pointer (*) or reference (&) is ignored.
;
; A method returns a TypeSpecifier type of data,
;
; Has a name (IDENTIFIER),
; and takes zero or more parameters.
;
; For languages that support 'const', it signifies that
; the named member is not allowed to modify the contents
; of the object, and has read-only access.
;
method: 'const'? TypeSpecifier '*'? IDENTIFIER '(' Param* (',' Param)* ')' 'const'? ';'

; An Param is of type Typespecifier, an optional pointer indicator '*',
; or an optional reference indicator '&', and an optional
; IDENTIFIER
.
Param : 'const'? TypeSpecifier ('*' | '&')? IDENTIFIER*

;
; A TypeSpecifier is any XDR data type
; OR any user defined data type.
;
TypeSpecifier : XdrDataType | UserDataType
\end{verbatim}
\caption{Class Method Definition}
\label{fig:ClassIntroXDR}
\end{figure}

\subsubsection{XDRDataType}
Where XdrDataType is defined
as \href {https://datatracker.ietf.org/doc/html/rfc4506#page-4}{XDR Data Types}

\subsubsection{IDENTIFIER}
An IDENTIFIER is a string starting with an upper or lower
case letter, followed by zero or more upper or lower case
letters or digits, or underline.

\begin{figure}
\begin{verbatim}
IDENTIFIER : [a-zA-Z] ([a-zA-Z0-9_])*
\end{verbatim}
\caption{IDENTIFIER Definition}
\label{fig:IDENTIFIERDefinitionXDR}
\end{figure}

\subsubsection{UserDataType}
A user data type, is any typedef, struct, enum, or class
name defined in the XDR file.
\begin{figure}
\begin{verbatim}
enum User1 { cat = 1, dog = 2, human = 3};

struct Car {
 uint16_t Year;
 string Make<>;
 string Model<>;
}

typedef struct Car Cars<>;
\end{verbatim}
\caption{UserDataType Definition}
\label{fig:UserDataTypeDefinitionXDR}
\end{figure}

In figure \ref{fig:UserDataTypeDefinitionXDR)}, User1, Car, and Cars
are examples of a UserDataType.


\section{ABNF Differences from RFC-5234}

\chapter{Command Line}
\section{Command Line}
RPCGEN++ can currently generate:
\begin{itemize}
\item C++ code from XDR (.x) files.
\item C++ from ABNF (.abnf) files.
\item ABNF files from XDR files.
\end{itemize}

\subsection{Command Line Arguments}
Here is a list of RPCGEN++ command line options.
Not all of these options are available at any time.
See the sections on generating code.

\begin{itemize}

\item --abnf

  Optional: When the input files is an XDR file,
  this produces ABNF out files for the XDR definitions.
  
\item --client

  Optional: Create a client that understands the input grammar
  and can interact with servers.
  
\item -Dname[=value]

  Optional: May occur more than once.
  The input files are processed with the c-preprocessor 'cpp'
  command. All '-D' values are passed to the c-preprocessor command.
  
\item --headers

  Optional: Produce header files needed by the lang.
  As other output languages are added, if they have definition
  files, this will create their definition files.
  
\item --input input-file
  
  Required: The name of the input file to be processed.
  Files that end with '.x' will be processed as XDR files.
  Files that end with '.abnf'' will be processed as ABNF files.
  
\item --lang=CPP
  
  Optional: Optional. If not supplied, CPP is assumed.
  More output languages will be added later.
  Currently only 'CPP' is supported.
  
\item --outdir out-directory-name
  
  Optional: If not specified the files will be generated
  in a directory named 'Generated'.
  
\item --server

  Optional: Create a server that understands the input grammar
  and responds to clients.
  
\item --stubs

  Optional: Create the source code wrappers for users to
  use as their staring point to get, set, transmit, and
  receive these objects over the network.
  
\item --xdr

  Optional: Produce the source code. These are the source
  code files that match the --headers option.
  
\item --xds

  Optional: Produce an XML XSD file.
  
\end{itemize}

\subsection{Generating C++ code from XDR files}

\subsection{Generating C++ code from ABNF files}

\subsection{Generating ABNF from XDR files}


\chapter{XDR files}
\section{Comments}
\subsection{Single Line Comments}
\subsection{Multiple Line Comments}

\section{Conditionals}

\section{Pass Through}

\section{Data Types}

\section{Fixed Size Arrays}

\section{Variable Sized Arrays}

\section{Strings}
\subsection{Fixed Size Strings}
\subsection{Variable Size Strings}
\subsection{Specifying a Maximum Size}

\section{Opaque Data}

\section{Member Functions}

\section{Bit Fields}

\chapter{ABNF files}
\chapter{Getting The source Code}
\section{Where to get RPCGEN++}

\chapter{Builing RPCGEN++ for yourself}
\section{Building Phoenix}
The entire Phoenix build system is designed to build everything
on Linux for all platforms using the \textit{make} command.

It uses the GNU compilers on Linux.
The mingw compilers to cross compile for Windows.
And the android NDK package to cross compile for Android.
The build has been tested on the Android based Oculus virtual
reality headsets.

The build uses the POSIX make command.
And can build debug, release, and test coverage libraries,
commands, and test code from the make system.

You can build any of the options below, or '\textbf{make all}
to build everything.

For each build, it builds 2 libraries and commands.
First using dynamic shared libraries, or crates them.
And the second pass builds with a minimum number
of dynamic libraries, using as many static libraries
as possible.

The BuildTools directory contains the Makefiles
that are included to do all of the different builds
and cross compiling.

You will need to modify the paths to the mingw and
android NDK paths to match your installation in the
Makefile.top and Makefile.top.lib files.

You will have to set the ANDROID\_CC, ANDROID\_CXX, ANDROID\_STRIP,
W64\_CC, W64\_CXX, and W64\_STRIP settings.

On Linux the mingw packages are often loaded using the operating
systems installation tools (dnf, apt, yum).

The Android NDK packages have to be downloaded using the
latest google instructions, which seem to change over time.
I used the NDK r26 version.

\subsection{Build Requirements}
You will need ... tools aasvg,
tex for documents ...
xdr2rfc for spects ...
g++ ...
gcov, gcovr, genhtml...
\subsection{Building For Linux}
From the top of the GIT repository you downloaded,
type one of these commands to build for Linux:

\subsubsection{Building debug for Linux}
Type \textbf{make LinuxDebug}

To build the Linux libraries and commands with symbol
and debugging information.
  
The output will be in BuildTemp/Linux/Debug.
The static library will be in BuildTemp/Linux/Debug/lib
and the shared library will be named lib{name}.so.{version}.
The static library will be named lib{name}.a.
  
Any commands built will be in BuildTemp/Linux/Debug/bin
for dynamic.

And BuildTemp/Linux/Debug/sbin for static programs.

\subsubsection{Building release for Linux}
Type \textbf{make LinuxRelease}

To build the Linux libraries and commands with the libraries
and programs compiled with the optimize switch turned on.
The necessary symbols will be stripped from them to make
them smaller.

The output will be in BuildTemp/Linux/Release.
The static library will be in BuildTemp/Linux/Release/lib
and the shared library will be named lib{name}.so.{version}.
The static library will be named lib{name}.a.
  
Any commands built will be in BuildTemp/Linux/Release/bin
for dynamic.

And BuildTemp/Linux/Release/sbin for static programs.

\subsubsection{Building test coverage for Linux}
Type \textbf{make LinuxTcov}

To build the Linux libraries and commands with the libraries
and programs compiled with the test code and test metric
collection information built into them.
See test coverage section \ref{sec:TestCoverage} for more information.
Debug symbols are included in TCOV builds.

The output will be in BuildTemp/Linux/Tcov.
The static library will be in BuildTemp/Linux/Tcov/lib
and the shared library will be named lib{name}.so.{version}.
The static library will be named lib{name}.a.

Any commands built will be in BuildTemp/Linux/Tcov/bin
for dynamic.

And BuildTemp/Linux/Tcov/sbin for static programs.

\textbf{When you run a TCOV build file or library, it generates a LOT of  output files}. Be sure to read section \ref{sec:TestCoverage}{Test Coverage}.
  
\subsubsection{Building everything for Linux}
Type \textbf{make Linux}

To build LinuxDebug, LinuxRelease, and LinuxTcov.

\subsection{Building For Windows on Linux}
  
From the top of the GIT repository you downloaded,
type one of these commands to build for Windows, 64-bit:


\subsubsection{Building debug for Windows}
Type \textbf{make W64Debug}

 To build the Windows 64-bit libraries and commands with symbol
 and debugging information.

 The output will be in BuildTemp/W64/Debug.
 The static library will be in BuildTemp/W64/Debug/lib
 and the shared library will be named lib{name}.so.{version}.
 The static library will be named lib{name}.a.
  
 Any commands built will be in BuildTemp/W65/Debug/bin,
 for dynamic.
  
 And BuildTemp/W64/Debug/sbin for static programs.

\subsubsection{Building releaes for Windows}
Type \textbf{make W64Release}

To build the Windows 64-bit libraries and commands with the libraries
and programs compiled with the optimize switch turned on.
The necessary symbols will be stripped from them to make
them smaller.

The output will be in BuildTemp/W64/Release.
The static library will be in BuildTemp/W64/Release/lib
and the shared library will be named lib{name}.so.{version}.
The static library will be named lib{name}.a.
  
Any commands built will be in BuildTemp/W65/Release/bin,
for dynamic.
  
And BuildTemp/W64/Release/sbin for static programs.
  
\subsubsection{Building test coverage for Windows}
Type \textbf{make W64Tcov}

To build the Windows libraries and commands with the libraries
and programs compiled with the test code and test metric
collection information built into them.
See test coverage section \ref{sec:TestCoverage} for more information.
Debug symbols are included in TCOV builds.

The output will be in BuildTemp/W64/Tcov.
The static library will be in BuildTemp/W64/Tcov/lib
and the shared library will be named lib{name}.so.{version}.
The static library will be named lib{name}.a.
  
Any commands built will be in BuildTemp/W64/Tcov/bin
for dynamic.

And BuildTemp/W64/Tcov/sbin for static programs.

\textbf{When you run a TCOV build file or library, it generates a LOT of  output files}. Be sure to read section \ref{sec:TestCoverage}{ Test Coverage}.
  
You have to run the built programs on a Widows device
with the matching GNU tools.
  
\subsubsection{Building everyting for Windows}
Type \textbf{make W64}

To build W64Debug, W64Release, and W64Tcov.

\subsection{Building For Android on Linux}
  
From the top of the GIT repository you downloaded,
type one of these commands to build for Android, 64-bit:
You will need the GCC NDK build tools in your path.

\subsubsection{Building debug for Android}
Type \textbf{make AndroidDebug}

To build the Android 64-bit libraries and commands with symbol
and debugging information.

The output will be in BuildTemp/Android/Debug.
The static library will be in BuildTemp/Android/Debug/lib
and the shared library will be named lib{name}.so.{version}.
The static library will be named lib{name}.a.
  
Any commands built will be in BuildTemp/Android/Debug/bin,
for dynamic.
  
And BuildTemp/Android/Debug/sbin for static programs.

\subsubsection{Building Android for Android}
Type \textbf{make AndroidRelease}

To build the Android 64-bit libraries and commands with the libraries
and programs compiled with the optimize switch turned on.
The necessary symbols will be stripped from them to make
them smaller.

The output will be in BuildTemp/Android/Release.
The static library will be in BuildTemp/W64/Release/lib
and the shared library will be named lib{name}.so.{version}.
The static library will be named lib{name}.a.
  
Any commands built will be in BuildTemp/Android/Release/bin,
for dynamic.
  
And BuildTemp/Android/Release/sbin for static programs.
  
\subsubsection{Building Test Coverage for Android}
Type \textbf{make AndroidTcov}

To build the Android libraries and commands with the libraries
and programs compiled with the test code and test metric
collection information built into them.
See test coverage section \ref{sec:TestCoverage} for more information.
Debug symbols are included in TCOV builds.

The output will be in BuildTemp/Android/Tcov.
The static library will be in BuildTemp/Android/Tcov/lib
and the shared library will be named lib{name}.so.{version}.
The static library will be named lib{name}.a.
  
Any commands built will be in BuildTemp/Android/Tcov/bin
for dynamic.

And BuildTemp/Android/Tcov/sbin for static programs.

\textbf{When you run a TCOV build file or library, it generates a LOT of  output files}. Be sure to read section \ref{sec:TestCoverage}{ Test Coverage}.

You have to run the built programs on a Android device
with the matching GNU tools.
  
\subsubsection{Building everything for Android}
Type \textbf{make Android}

To build AndroidDebug, AndroidRelease, and AndroidTcov.

\subsection{Building the Documentation}
The documentation for RPCGEN++ is built using latex and make4ht.
You will need those.
Go to the NewRpcgen/Documentation/LaTex directory
and type 'make'.

\subsection{Filing Bugs Reports}
File bug, issues, and questions on github, at:
\href {https://github.com/RiverExplorer/Phoenix/issues}{Bug and Issue Reports}

\section{Test Coverage}
This project uses \textit{gcov} and \textit{gcovr} to generate the HTML
test coverage results.

This produces a lot of data and this section will get you
started in understanting what you have produced.

Once the program or library is compiled with the test coverage
tools, you must run them on the target operating system.
So the Windows 64-bit builds (W64) must be run on a windows
computer to get the test results for windows.
And same with the Android devices.

More ....


\chapter{Legal}
\section{Legal Mumbo Jumbo}
\centering{
  \textbf{This document is Copyright 2025 Douglas Mark Royer}.
  }

  \begin{figure}
    \centering
    \includegraphics{RiverExplorerLLC-Logo-BLACK-131x50-Transparent.ps.png}
    \caption{RiverExplorer Company Logo}
    \label{fig:RiverExplorerLogo}
  \end{figure}
  
\centering {
  \textbf{'RiverExplorer' and the logo shown in figure \ref{fig:RiverExplorerLogo} are trademarks of Douglas Mark Royer}.
  \textbf{'RiverExplorer Games LLC' is a wholly owned subsidiary of RiverExplorer LLC.}
  }
\subsection{License}
\textit{Permission is hereby granted, free of charge, to any person obtaining a copy of this software and associated documentation files (the “Software”), to deal in the Software without restriction, including without limitation the rights to use, copy, modify, merge, publish, distribute, sublicense, and/or sell copies of the Software, and to permit persons to whom the Software is furnished to do so, subject to the following conditions:}

\textit{The above copyright notice and this permission notice shall be included in all copies or substantial portions of the Software.}

\textit{THE SOFTWARE IS PROVIDED “AS IS”, WITHOUT WARRANTY OF ANY KIND, EXPRESS OR IMPLIED, INCLUDING BUT NOT LIMITED TO THE WARRANTIES OF MERCHANTABILITY, FITNESS FOR A PARTICULAR PURPOSE AND NONINFRINGEMENT. IN NO EVENT SHALL THE AUTHORS OR COPYRIGHT HOLDERS BE LIABLE FOR ANY CLAIM, DAMAGES OR OTHER LIABILITY, WHETHER IN AN ACTION OF CONTRACT, TORT OR OTHERWISE, ARISING FROM, OUT OF OR IN CONNECTION WITH THE SOFTWARE OR THE USE OR OTHER DEALINGS IN THE SOFTWARE.}

\appendix
\chapter{Sample XDR files}
\chapter{Sample ABNF files}
\end{document}
